\section{Results}

\subsection{Feature Coverage}
Table~\ref{tab:features} summarizes the proposed features from the original project goals alongside their implementation status.

\begin{table}[htbp]
\caption{Proposed vs.\ Implemented Features}
\label{tab:features}
\centering
\renewcommand{\arraystretch}{1.3}
\rowcolors{2}{white}{gray!10}
\begin{tabularx}{\columnwidth}{|X|c|}
\hline
\rowcolor[HTML]{4472C4}
\textcolor{white}{\textbf{Feature}} & \textcolor{white}{\textbf{Status}} \\ \hline
CLI accepting \texttt{.ll} files               & \checkmark \\ \hline
LLVM IR parsing via \texttt{llvmlite}          & \checkmark \\ \hline
Typed CFG edge extraction                       & \checkmark \\ \hline
Intermediate scene graph                        & \checkmark \\ \hline
JSON scene graph export                         & \checkmark \\ \hline
Graphviz DOT/PNG export                         & \checkmark \\ \hline
Basic stack animation                           & \checkmark \\ \hline
Rich IR spotlight animation                     & \checkmark \\ \hline
Rich-SSA 3-column animation                     & \checkmark \\ \hline
CFG traversal animation                         & \checkmark \\ \hline
Runtime trace import/export                     & \checkmark \\ \hline
Domtree/loop metadata import                    & \checkmark \\ \hline
GIF output with palette workflow                & \checkmark \\ \hline
Adjustable animation speed                      & \checkmark \\ \hline
Automatic C compilation                         & $\times$   \\ \hline
Interactive stepping/pause                      & $\times$   \\ \hline
YAML configuration files                        & $\times$   \\ \hline
Color palette selection                         & $\times$   \\ \hline
GUI                                             & $\times$   \\ \hline
\end{tabularx}
\end{table}

Of the originally proposed proof-of-concept targets, all were achieved except automatic C-to-IR compilation (users supply \texttt{.ll} files directly) and interactive stepping. The three-column SSA panel and CFG traversal animation were implemented as additional capabilities beyond the original scope.

\subsection{Supported IR Instructions}
LLVManim classifies each LLVM IR instruction into one of nine \texttt{EventKind} categories:

\begin{itemize}
    \item \texttt{alloca}: Stack memory allocation
    \item \texttt{load} / \texttt{store}: Memory read and write operations
    \item \texttt{binop}: Binary arithmetic operations (\texttt{add}, \texttt{sub}, \texttt{mul}, etc.)
    \item \texttt{compare}: Integer and floating-point comparisons (\texttt{icmp}, \texttt{fcmp})
    \item \texttt{call}: Function invocations
    \item \texttt{ret}: Function returns
    \item \texttt{br}: Branch instructions (conditional and unconditional)
    \item \texttt{other}: Unsupported or unrecognized instructions (skipped with a diagnostic message)
\end{itemize}

CFG edge extraction supports five terminator types: \texttt{br}, \texttt{switch}, \texttt{invoke}, \texttt{indirectbr}, and \texttt{callbr}. Conditional branches are labeled with T/F direction annotations.

\subsection{Animation Modes}
LLVManim provides three stack-focused animation modes and one graph-based mode. Each mode can be reproduced with a single CLI command.

\textbf{Basic mode} (\texttt{--ir-mode basic}) renders a stack-only view where each instruction triggers a per-cell badge flash:
\begin{lstlisting}[language=bash,gobble=0]
uv run llvmanim input.ll --animate \
    --ir-mode basic --outdir out/
\end{lstlisting}

\textbf{Rich mode} (\texttt{--ir-mode rich}) adds a two-column layout with an IR source panel and a moving spotlight cursor highlighting the currently executing instruction:
\begin{lstlisting}[language=bash,gobble=0]
uv run llvmanim input.ll --animate \
    --ir-mode rich --outdir out/
\end{lstlisting}

\textbf{Rich-SSA mode} (\texttt{--ir-mode rich-ssa}) extends the rich layout to three columns (IR Source $\mid$ SSA Values $\mid$ Stack), displaying computed symbolic values for binary operations, comparisons, and loads:
\begin{lstlisting}[language=bash,gobble=0]
uv run llvmanim input.ll --animate \
    --ir-mode rich-ssa --outdir out/
\end{lstlisting}

\textbf{CFG traversal} (\texttt{--cfg-animate}) renders an animated walk through the control-flow graph, using Graphviz DOT layout for node positions and a runtime trace for the traversal order:
\begin{lstlisting}[language=bash,gobble=0]
uv run llvmanim input.ll --cfg-animate \
    --dot-cfg main.dot \
    --import-trace trace.json --outdir out/
\end{lstlisting}

% Placeholder: add screenshots of each mode here when available.
% \begin{figure*}[htbp]
% \centerline{\includegraphics[width=\textwidth]{basic_mode.png}}
% \caption{Basic mode animation showing stack operations for \texttt{double.ll}.}
% \label{fig:basic_mode}
% \end{figure*}

\subsection{Test Coverage and Quality}
The LLVManim implementation is validated by 302 automated tests organized into four marker categories:

\begin{table}[htbp]
\caption{Test Distribution by Category}
\label{tab:tests}
\centering
\renewcommand{\arraystretch}{1.3}
\rowcolors{2}{white}{gray!10}
\begin{tabularx}{\columnwidth}{|X|c|}
\hline
\rowcolor[HTML]{4472C4}
\textcolor{white}{\textbf{Category}} & \textcolor{white}{\textbf{Count}} \\ \hline
Unit (ingest, transform)        & 192 \\ \hline
Integration (present, cli, pipeline) & 96 \\ \hline
Contract (export boundaries)    & 8   \\ \hline
End-to-end (entrypoints)        & 6   \\ \hline
\textbf{Total}                  & \textbf{302} \\ \hline
\end{tabularx}
\end{table}

Line coverage stands at 84\% across the entire \texttt{llvmanim} package. A four-stage quality gate runs on every pull request:
\begin{enumerate}
    \item \textbf{Ruff}: Linting and import sorting
    \item \textbf{Pyright}: Static type checking (strict mode)
    \item \textbf{import-linter}: Architectural boundary enforcement between layers
    \item \textbf{pytest}: Full test suite execution with coverage reporting
\end{enumerate}

\subsection{Reproducibility}
All LLVManim results can be reproduced from the public repository\cite{llvmanim-repo}:

\begin{lstlisting}[language=bash,gobble=0]
git clone https://github.com/BenHoule/LLVManim.git
cd LLVManim
uv sync --dev
./scripts/quality-check.sh  # runs all 4 gates
\end{lstlisting}

System requirements: Python 3.12+, \texttt{uv}, Linux packages \texttt{libcairo2-dev} and \texttt{libpango1.0-dev} (for Manim), and optionally \texttt{graphviz} for PNG graph export.
